\documentclass[11pt,a4paper]{article}
\usepackage[italian]{babel}
\usepackage[margin=2.5cm]{geometry}
\usepackage{parskip}
\usepackage{amsmath, amssymb, amsthm}
\usepackage{booktabs}
\usepackage{array}
\usepackage{xcolor}
\usepackage{needspace}
\usepackage{enumitem}
\usepackage{listings}
\usepackage{tikz}
\usepackage[most]{tcolorbox}
\usepackage[hidelinks]{hyperref}
\usetikzlibrary{decorations.pathreplacing, shapes.geometric, arrows.meta, positioning}

% Niente vedove/orfane: i paragrafi non lasciano righe isolate a fine/inizio pagina
\widowpenalty=10000
\clubpenalty=10000

% Elenchi più compatti
\setlist{itemsep=2pt, parsep=0pt, topsep=4pt, partopsep=0pt}

% --- Riquadri con bordo nero, senza numero (si spezzano tra due pagine) ---
% Uso: \begin{definizione} ... \end{definizione}
%      \begin{teorema}[Nome] ... \end{teorema}  (il nome è facoltativo)
\tcbset{nota/.style={enhanced jigsaw, breakable, colback=white,
  colframe=black, boxrule=0.5pt, sharp corners}}
\NewTColorBox{definizione}{}{nota,
  before upper={\textbf{Definizione.}\ }}
\NewTColorBox{teorema}{o}{nota,
  before upper={\textbf{Teorema\IfValueT{#1}{ (#1)}.}\ }}
\NewTColorBox{proposizione}{o}{nota,
  before upper={\textbf{Proposizione\IfValueT{#1}{ (#1)}.}\ }}
\NewTColorBox{proprieta}{o}{nota,
  before upper={\textbf{Proprietà\IfValueT{#1}{ (#1)}.}\ }}
\NewTColorBox{assioma}{o}{nota, boxrule=1pt,
  before upper={\textbf{Assioma\IfValueT{#1}{ (#1)}.}\ }}

% --- Ambienti semplici (senza riquadro) ---
% Il titolo sta in cima al blocco, anche se il contenuto inizia con un riquadro o
% con codice; e non resta mai da solo in fondo a una pagina.
% Uso: \begin{esempio}[Titolo facoltativo] ... \end{esempio}
\newcommand{\newrunin}[2]{%
  \NewDocumentEnvironment{#1}{o}{%
    \par\Needspace{4\baselineskip}\addvspace{\medskipamount}\noindent
    \textbf{#2}\IfValueT{##1}{ (##1)}\textbf{.}\ \ignorespaces}%
  {\par\addvspace{\medskipamount}}}
\newrunin{esempio}{Esempio}
\newrunin{esercizio}{Esercizio}
\newrunin{osservazione}{Osservazione}

% V e F nelle tavole di verità
\newcommand{\V}{\textsf{V}}
\newcommand{\F}{\textsf{F}}

% --- Codice C ---
% Uso: \begin{codice} ... \end{codice}      codice C numerato
%      \begin{comando} ... \end{comando}    comandi da terminale
%      \begin{risultato} ... \end{risultato} output di un programma
%      \lstinline|printf("ciao");|           codice dentro al testo
\lstdefinestyle{c}{
  language=C,
  basicstyle=\ttfamily\small,
  keywordstyle=\color{blue}\bfseries,
  commentstyle=\color{gray}\itshape,
  stringstyle=\color{red!70!black},
  numbers=left,
  numberstyle=\tiny\color{gray},
  numbersep=8pt,
  columns=flexible,
  keepspaces=true,
  breaklines=true,
  showstringspaces=false,
  tabsize=4,
  morekeywords={include, define},
  % lettere accentate nei commenti
  literate={à}{{\`a}}1 {è}{{\`e}}1 {é}{{\'e}}1 {ì}{{\`i}}1
           {ò}{{\`o}}1 {ù}{{\`u}}1 {È}{{\`E}}1 {À}{{\`A}}1
}
\lstset{style=c}
\newtcblisting{codice}{listing only, nota, left=6mm,
  listing options={style=c}}
\newtcblisting{comando}{listing only, nota,
  listing options={basicstyle=\ttfamily\small, numbers=none, language={},
    columns=flexible, keepspaces=true, breaklines=true}}
\newtcblisting{risultato}{listing only, enhanced jigsaw, breakable,
  colback=black!5, colframe=black, boxrule=0.5pt, sharp corners,
  listing options={basicstyle=\ttfamily\small, numbers=none, language={},
    columns=flexible, keepspaces=true, breaklines=true}}

% --- Diagrammi di flusso (simboli standard) ---
\tikzset{
  terminale/.style = {ellipse, draw, minimum width=2.5cm, minimum height=0.9cm},
  io/.style        = {trapezium, trapezium left angle=70, trapezium right angle=110,
                      draw, minimum width=2.5cm, minimum height=0.9cm},
  processo/.style  = {rectangle, draw, minimum width=2.5cm, minimum height=0.9cm},
  decisione/.style = {diamond, draw, aspect=2, minimum width=2.5cm},
  freccia/.style   = {thick, -{Stealth}}
}

\title{Gruppi}
\author{Marco Cerbella}
\date{Lezione 9 -- 19 ottobre 2026}

\begin{document}
\maketitle

\section{Gruppi}

\begin{definizione}
$(G, *)$ è un \emph{gruppo} se
\begin{enumerate}
  \item $*$ è associativa: $\forall a, b, c \in G$ \quad $(a * b) * c = a * (b * c)$;
  \item esiste l'elemento neutro: $\exists e \in G$ tale che $\forall a \in G$ \quad $e * a = a * e = a$;
  \item ogni elemento è invertibile: $\forall a \in G$ \quad $\exists a' \in G$ tale che $a * a' = a' * a = e$.
\end{enumerate}
Un gruppo $(G, *)$ è \emph{abeliano} se $*$ è commutativa: $\forall a, b \in G$ \quad $a * b = b * a$.
\end{definizione}

\begin{esempio}
\begin{itemize}
  \item $(\mathbb{Q}\setminus\{0\}, \cdot)$ è un gruppo commutativo: l'operazione è associativa, l'elemento neutro è $1$ e l'inverso di $\frac{a}{b}$ è $\frac{b}{a}$ (se $q \in \mathbb{Q}$ l'inverso è $q^{-1}$).
  \item $(\mathbb{Z}, +)$ è un gruppo abeliano: associativa, $0$ elemento neutro, e se $k \in \mathbb{Z}$ allora $-k$ è l'inverso.
\end{itemize}
\end{esempio}

\subsection{Notazione}

\begin{center}
\begin{tabular}{@{}p{3.5cm} p{4.5cm} p{4.5cm}@{}}
\toprule
$(G, *)$ gruppo & Notazione additiva $(G, +)$ & Notazione moltiplicativa $(G, \cdot)$ \\
\midrule
$e$ elemento neutro & $0$ elemento neutro & $1$ elemento neutro \\[2pt]
$a'$ inverso di $a$ & $(-a)$ \emph{opposto} di $a$ & $a^{-1}$ \emph{inverso} di $a$ \\[2pt]
$\underbrace{a * a * \dots * a}_{n\text{ volte}}$ & $\underbrace{a + \dots + a}_{n\text{ volte}} = na$ & $\underbrace{a \cdot a \cdot \ldots \cdot a}_{n\text{ volte}} = a^n$ \\
\bottomrule
\end{tabular}
\end{center}

\begin{definizione}[Gruppo ciclico]
Un gruppo $(G, *)$ è un \emph{gruppo ciclico} se esiste un elemento $g \in G$ tale che $\forall a \in G$
\[
  a = \underbrace{g * g * \dots * g}_{\text{qualche volta}}
  \quad\text{oppure}\quad a = e\ (= g^0)
  \quad\text{oppure}\quad a = g^{-1} * g^{-1} * \dots * g^{-1} .
\]
\end{definizione}

\begin{esempio}[$\mathbb{Z}_{12}$: l'aritmetica dell'orologio]
$\mathbb{Z}_{12} = \{[1], [2], \dots, [12]\}$ con la somma ad esempio $[3] + [4] = [7]$, $[8] + [5] = [1]$. Si può costruire la tavola della somma, con $[12]$ elemento neutro e ogni elemento invertibile (ad esempio $[5] + [?] = [12]$ dà $[?] = [7]$). Poiché $[2] = [1] + [1]$, $[3] = [1] + [1] + [1]$, \dots, il gruppo è ciclico, generato da $[1]$.
\end{esempio}

\begin{esempio}[$(\mathbb{Z}, +)$ è ciclico]
$(\mathbb{Z}, +)$ è un gruppo ciclico: $+1$ genera tutti i numeri. Infatti, se $h \in \mathbb{Z}$:
\begin{itemize}
  \item $h > 0$: $h = \underbrace{1 + 1 + \dots + 1}_{h\text{ volte}}$;
  \item $h = 0$: elemento neutro;
  \item $h < 0$: $h = \underbrace{(-1) + (-1) + \dots + (-1)}_{(-h)\text{ volte}}$.
\end{itemize}
\end{esempio}

\subsection{Simmetrie di un rettangolo}

\begin{esempio}
Consideriamo un rettangolo $ABCD$ (con $A$, $B$ in basso e $D$, $C$ in alto) e le sue simmetrie:
\begin{itemize}
  \item $a_1$: simmetria rispetto all'asse verticale, $(ABCD) \mapsto (BADC)$;
  \item $a_2$: simmetria rispetto all'asse orizzontale, $(ABCD) \mapsto (DCBA)$;
  \item $\tau$: rotazione di $180^\circ$, $(ABCD) \mapsto (CDAB)$;
  \item $\mathrm{Id}$: identità, $(ABCD) \mapsto (ABCD)$.
\end{itemize}
La composizione di simmetrie è ancora una simmetria del rettangolo. Ad esempio
\[
  \mathrm{Id} * \tau\,(ABCD) = \mathrm{Id}(CDAB) = (CDAB), \qquad
  \tau * \tau\,(ABCD) = \tau(CDAB) = (ABCD),
\]
\[
  \tau * a_1\,(ABCD) = \tau(BADC) = (DCBA) = a_2\,(ABCD).
\]
Dunque $\tau * a_1 = a_2$ (il resto della tavola della composizione viene completato nella lezione successiva).
\end{esempio}

\end{document}
